\newpage
\begin{section}{Introduction}
    Indoor positioning systems have various use cases for industries such as tracking inventory or navigating large buildings \cite{Wang2026}.
    However, existing solutions often struggle in these environments: GPS suffers from indoor signal attenuation and multipath interference \cite{Sysnav2024}, while NFC is limited by short range \cite{Workm82024}.
    Other systems like camera-based setups limit privacy \cite{Chelani2025} and Bluetooth devices like AirTags lack real-time updates \cite{VandeVelde2024}. 
    Thus, the requirements for an effective indoor positioning system include low interference and cost and no line-of-sight requirements.
    We therefore, investigate an electromagnetic positioning system designed specifically for 2D indoor localization, using model trains as the objects to be tracked.

    \begin{subsection}{Problem statement}
        Electromagnetic solutions have their own engineering challenges such as distortion caused by picking up  electromagnetic interference.
        Furthermore, the system's range is limited by the strength of its magnetic field, which diminishes with distance. 
        Compensating for this requires either increasing transmitter power or using larger coils. 
        Because power draw is constrained by the train setup and physical dimensions are limited by the scale of the locomotive, achieving sufficient operating range presents a significant design challenge.
    \end{subsection}
    
    \begin{subsection}{Scope/Limitations}
        data aquisition using sound cardsinstead of ?
    \end{subsection}
  
    \begin{subsection}{Report structure}
        The remainder of this report is organized as follows: Section 3 presents related work used for research, Section 4 gives an overview of the system, Section 5 describes the power supply for the system, Section 6 describes the transmitter, Section 7 describes the receiver, Section 8 describes the positioning algorithm, Section 9 discusses the overall resulting system and Section 10 concludes the report.
        Section 5, 6, 7, 8 are further structured into method, results and discussion subsections. 
    \end{subsection}
    
\end{section}